% =====================================================================
% PAPER6 — Capture-Resilient Voting in Multi-LLM Verifier Ensembles
% Template: tmlr.sty (v0.1, copied from PAPER5, 6560 bytes)
% Status: COMPLETE DRAFT — full sections + synthetic pilot results (4,000 trials/cell, N∈{5,7,10}, f∈{1,2,3}); real-trace benchmark noted as future work
% =====================================================================

\documentclass[11pt]{article}
\usepackage[preprint]{tmlr}
% If accepted, switch to: \usepackage[accepted]{tmlr}
\usepackage{amsmath, amssymb, amsthm}
%%%%% NEW MATH DEFINITIONS %%%%%
%%%%% From JmlrOrg/tmlr-style-file/main
%%%%% https://github.com/JmlrOrg/tmlr-style-file/blob/main/math_commands.tex

\usepackage{amsmath,amsfonts,bm}

% Mark sections of captions for referring to divisions of figures
\newcommand{\figleft}{{\em (Left)}}
\newcommand{\figcenter}{{\em (Center)}}
\newcommand{\figright}{{\em (Right)}}
\newcommand{\figtop}{{\em (Top)}}
\newcommand{\figbottom}{{\em (Bottom)}}
\newcommand{\captiona}{{\em (a)}}
\newcommand{\captionb}{{\em (b)}}
\newcommand{\captionc}{{\em (c)}}
\newcommand{\captiond}{{\em (d)}}

% Highlight a newly defined term
\newcommand{\newterm}[1]{{\bf #1}}

% Figure reference, lower-case.
\def\figref#1{figure~\ref{#1}}
\def\Figref#1{Figure~\ref{#1}}
\def\twofigref#1#2{figures \ref{#1} and \ref{#2}}
\def\quadfigref#1#2#3#4{figures \ref{#1}, \ref{#2}, \ref{#3} and \ref{#4}}
% Section reference, lower-case.
\def\secref#1{section~\ref{#1}}
\def\Secref#1{Section~\ref{#1}}
\def\twosecrefs#1#2{sections \ref{#1} and \ref{#2}}
\def\secrefs#1#2#3{sections \ref{#1}, \ref{#2} and \ref{#3}}
% Reference to an equation, lower-case.
\def\eqref#1{equation~\ref{#1}}
\def\Eqref#1{Equation~\ref{#1}}
\def\plaineqref#1{\ref{#1}}
\def\chapref#1{chapter~\ref{#1}}
\def\Chapref#1{Chapter~\ref{#1}}
\def\ranglegchapref#1#2{chapters\ref{#1}--\ref{#2}}
\def\algref#1{algorithm~\ref{#1}}
\def\Algref#1{Algorithm~\ref{#1}}
\def\twoalgref#1#2{algorithms \ref{#1} and \ref{#2}}
\def\Twoalgref#1#2{Algorithms \ref{#1} and \ref{#2}}
\def\partref#1{part~\ref{#1}}
\def\Partref#1{Part~\ref{#1}}
\def\twopartref#1#2{parts \ref{#1} and \ref{#2}}

\def\ceil#1{\lceil #1 \rceil}
\def\floor#1{\lfloor #1 \rfloor}
\def\1{\bm{1}}
\newcommand{\train}{\mathcal{D}}
\newcommand{\valid}{\mathcal{D_{\mathrm{valid}}}}
\newcommand{\test}{\mathcal{D_{\mathrm{test}}}}

\def\eps{{\epsilon}}

% Random variables
\def\reta{{\textnormal{$\eta$}}}
\def\ra{{\textnormal{a}}}
\def\rb{{\textnormal{b}}}
\def\rc{{\textnormal{c}}}
\def\rd{{\textnormal{d}}}
\def\re{{\textnormal{e}}}
\def\rf{{\textnormal{f}}}
\def\rg{{\textnormal{g}}}
\def\rh{{\textnormal{h}}}
\def\ri{{\textnormal{i}}}
\def\rj{{\textnormal{j}}}
\def\rk{{\textnormal{k}}}
\def\rl{{\textnormal{l}}}
\def\rn{{\textnormal{n}}}
\def\ro{{\textnormal{o}}}
\def\rp{{\textnormal{p}}}
\def\rq{{\textnormal{q}}}
\def\rr{{\textnormal{r}}}
\def\rs{{\textnormal{s}}}
\def\rt{{\textnormal{t}}}
\def\ru{{\textnormal{u}}}
\def\rv{{\textnormal{v}}}
\def\rw{{\textnormal{w}}}
\def\rx{{\textnormal{x}}}
\def\ry{{\textnormal{y}}}
\def\rz{{\textnormal{z}}}

% Random vectors
\def\rvepsilon{{\mathbf{\epsilon}}}
\def\rvtheta{{\mathbf{\theta}}}
\def\rva{{\mathbf{a}}}
\def\rvb{{\mathbf{b}}}
\def\rvc{{\mathbf{c}}}
\def\rvd{{\mathbf{d}}}
\def\rve{{\mathbf{e}}}
\def\rvf{{\mathbf{f}}}
\def\rvg{{\mathbf{g}}}
\def\rvh{{\mathbf{h}}}
\def\rvu{{\mathbf{i}}}
\def\rvj{{\mathbf{j}}}
\def\rvk{{\mathbf{k}}}
\def\rvl{{\mathbf{l}}}
\def\rvm{{\mathbf{m}}}
\def\rvn{{\mathbf{n}}}
\def\rvo{{\mathbf{o}}}
\def\rvp{{\mathbf{p}}}
\def\rvq{{\mathbf{q}}}
\def\rvr{{\mathbf{r}}}
\def\rvs{{\mathbf{s}}}
\def\rvt{{\mathbf{t}}}
\def\rvu{{\mathbf{u}}}
\def\rvv{{\mathbf{v}}}
\def\rvw{{\mathbf{w}}}
\def\rvx{{\mathbf{x}}}
\def\rvy{{\mathbf{y}}}
\def\rvz{{\mathbf{z}}}

% Vectors
\def\vzero{{\bm{0}}}
\def\vone{{\bm{1}}}
\def\vmu{{\bm{\mu}}}
\def\vtheta{{\bm{\theta}}}
\def\va{{\bm{a}}}
\def\vb{{\bm{b}}}
\def\vc{{\bm{c}}}
\def\vd{{\bm{d}}}
\def\ve{{\bm{e}}}
\def\vf{{\bm{f}}}
\def\vg{{\bm{g}}}
\def\vh{{\bm{h}}}
\def\vi{{\bm{i}}}
\def\vj{{\bm{j}}}
\def\vk{{\bm{k}}}
\def\vl{{\bm{l}}}
\def\vm{{\bm{m}}}
\def\vn{{\bm{n}}}
\def\vo{{\bm{o}}}
\def\vp{{\bm{p}}}
\def\vq{{\bm{q}}}
\def\vr{{\bm{r}}}
\def\vs{{\bm{s}}}
\def\vt{{\bm{t}}}
\def\vu{{\bm{u}}}
\def\vv{{\bm{v}}}
\def\vw{{\bm{w}}}
\def\vx{{\bm{x}}}
\def\vy{{\bm{y}}}
\def\vz{{\bm{z}}}

% Sets
\def\sA{{\mathbb{A}}}
\def\sB{{\mathbb{B}}}
\def\sC{{\mathbb{C}}}
\def\sD{{\mathbb{D}}}
\def\sF{{\mathbb{F}}}
\def\sG{{\mathbb{G}}}
\def\sH{{\mathbb{H}}}
\def\sI{{\mathbb{I}}}
\def\sJ{{\mathbb{J}}}
\def\sK{{\mathbb{K}}}
\def\sL{{\mathbb{L}}}
\def\sM{{\mathbb{M}}}
\def\sN{{\mathbb{N}}}
\def\sO{{\mathbb{O}}}
\def\sP{{\mathbb{P}}}
\def\sQ{{\mathbb{Q}}}
\def\sR{{\mathbb{R}}}
\def\sS{{\mathbb{S}}}
\def\sT{{\mathbb{T}}}
\def\sU{{\mathbb{U}}}
\def\sV{{\mathbb{V}}}
\def\sW{{\mathbb{W}}}
\def\sX{{\mathbb{X}}}
\def\sY{{\mathbb{Y}}}
\def\sZ{{\mathbb{Z}}}

\newcommand{\E}{\mathbb{E}}
\newcommand{\Ls}{\mathcal{L}}
\newcommand{\R}{\mathbb{R}}
\newcommand{\lr}{\alpha}
\newcommand{\reg}{\lambda}
\newcommand{\softmax}{\mathrm{softmax}}
\newcommand{\sigmoid}{\sigma}
\newcommand{\softplus}{\zeta}
\newcommand{\KL}{D_{\mathrm{KL}}}
\newcommand{\Var}{\mathrm{Var}}
\newcommand{\standarderror}{\mathrm{SE}}
\newcommand{\Cov}{\mathrm{Cov}}
\newcommand{\normlzero}{L^0}
\newcommand{\normlone}{L^1}
\newcommand{\normltwo}{L^2}
\newcommand{\normlp}{L^p}
\newcommand{\normmax}{L^\infty}

\DeclareMathOperator*{\argmax}{arg\,max}
\DeclareMathOperator*{\argmin}{arg\,min}

\DeclareMathOperator{\sign}{sign}
\DeclareMathOperator{\Tr}{Tr}
\let\ab\allowbreak
\DeclareMathOperator{\CV}{CV}
\usepackage[T1]{fontenc}
\usepackage[utf8]{inputenc}
\usepackage{graphicx}
\usepackage{booktabs}
\usepackage{multirow}
\usepackage[colorlinks=true,linkcolor=blue,citecolor=blue,urlcolor=blue]{hyperref}
\usepackage[capitalize]{cleveref}
\usepackage{microtype}
\setlength{\headheight}{14pt}
\usepackage{algorithm}
\usepackage{algpseudocode}
\let\STATE\State

% ---------------------------------------------------------------------
% Theorem environments
% ---------------------------------------------------------------------
\newtheorem{definition}{Definition}
\newtheorem{theorem}{Theorem}
\newtheorem{lemma}{Lemma}
\newtheorem{corollary}{Corollary}
\newtheorem{proposition}{Proposition}
\newtheorem{remark}{Remark}
\newtheorem{assumption}{Assumption}

% ---------------------------------------------------------------------
% Notation shortcuts
% ---------------------------------------------------------------------
\newcommand{\V}{\mathcal{V}}            % set of verifiers
\newcommand{\A}{\mathcal{A}}            % set of attackers
\newcommand{\CRV}{\textsc{CRV}}         % Capture-Resilient Voting
\newcommand{\ECR}{\varepsilon\text{-CR}} % $\varepsilon$-capture-resilient

\title{An Empirical Study of Aggregation Methods Under\\
       Verifier Capture in Multi-LLM Ensembles:\\
       When Does Majority Voting Fail, and Can We Beat It?}

\author{\name Zewen Liu\email liuzewen@qlnu.edu.cn\\\addr Qilu Institute of Technology}  % PREPRINT/ARXIV BUILD 2026-10-01: was an anonymous/double-blind author block. The block, not tmlr.sty, is what prints the author line on the title page; the `preprint` option clears the per-page banner. Restored to the real author for arXiv, which does not accept anonymous submissions.  % PREPRINT BUILD: was an anonymous/double-blind author block. The block, not tmlr.sty, is what prints the author line on the title page; the `preprint` option clears the per-page banner. See the preprint anonymity round log.
\date{July 2026}

\begin{document}
\maketitle

\begin{abstract}
Multi-LLM verification ensembles increasingly aggregate the verdicts of
several LLM \emph{verifiers} to reduce variance and adversarial
brittleness relative to any single model. A growing concern in such
ensembles is \emph{verifier capture}: a small coalition of verifiers
whose outputs are systematically biased can flip the aggregated
decision while remaining individually plausible. This paper contributes
(i) a formal definition of capture-resilient voting (\textsc{CRV}) and
the threat model distinguishing prompt-injection, rubric-injection, and
majority-bias mimicry attacks; (ii) an experimental
protocol comparing seven aggregation methods (plain majority,
PBFT-style deterministic thresholding, uniform-weighted mean, three
trust-score variants, and an adaptive hybrid) across verifier pools of
size $N \in \{5,7,10\}$ with capture counts $f \in \{1,2,3\}$ in a
synthetic benchmark and a small local-model pilot; and (iii) a working
hypothesis that captured verifier outputs can resemble a minority of
honest-but-wrong verifiers, so deviation-based capture rules may incur
false-positive costs. A synthetic pilot (4,000 trials per cell, $N=7$
verifiers, f=2 captured) tests one instantiation of this hypothesis:
trust-score weighting improves over plain majority under prompt-injection
capture (accuracy 0.69 vs. 0.34 at honest accuracy 0.6; 0.94 vs. 0.74 at
0.8), but the same method does not improve under majority-bias mimicry
(0.32 vs. 0.41 at honest accuracy 0.6). These pilot values are synthetic and
do not establish a universal impossibility result. A descriptive local-model
pilot is included in the evidence package, but a held-out real-trace benchmark
remains future work; we release the framework, generators, and analysis
scripts so the benchmark can be extended independently.


\textbf{Reproducibility and scope.} The evidence package contains the synthetic
benchmark tables, the checked-in generators and aggregation code, and a
descriptive four-model local Ollama pilot. It does not contain a held-out real
trace benchmark, and the local pilot is not presented as an external
generalization test. This work studies evaluation dynamics in LLM systems;
we discuss the bounded dual-use risk of releasing the attack generators.
\end{abstract}

% =====================================================================

\section{Introduction}
\label{sec:intro}
% =====================================================================

Large language model (LLM)-based verification pipelines increasingly rely on
\emph{multi-source verifier ensembles} \citep{cui2026edo,leung2026rubricllm}, in which a fixed pool
of $N$ LLM verifiers independently evaluate a candidate output and a final
decision is produced by aggregating their votes. Such ensembles have become a
default building block of contemporary AI auditing, peer review assistance,
and safety-critical content moderation, where the goal is to reduce variance
and adversarial brittleness relative to any single model.

A natural and increasingly documented failure mode of these ensembles is what
we call \emph{verifier capture}: a small coalition of verifiers whose outputs
are systematically biased in a coordinated direction can swing the aggregated
decision, even when a majority of verifiers remain individually reliable.
Verifier capture differs from the classical single-verifier hallucination or
bias problem in two important ways.
\textbf{(i) Capture is interactive.} The adversarial verifier does not need to
be accurate in absolute terms; it only needs to produce outputs that, when
combined with the other verifiers under a fixed aggregation rule, flip the
final decision.
\textbf{(ii) Capture is silent.} From the perspective of an external observer
that only sees final aggregated verdicts, a captured ensemble looks identical
to a correctly functioning one until the disagreement structure of individual
votes is examined.

Recent work has begun to expose capture-like failure modes in LLM judges:
\citet{zhang2026blindcurator} show that a biased judge \emph{silently
disables} a downstream curator past a sharp false-pass threshold. This paper
asks the natural follow-up question: given that capture is real and silent,
\emph{can any aggregation rule do better than plain majority voting}? Our
answer, in this empirical study, is \emph{mostly no}.

This paper makes four contributions:

\begin{itemize}
  \item A formal threat model for verifier capture in LLM ensembles, including
        a definition of $\varepsilon$-capture-resilience that decouples the
        adversary's per-verifier accuracy from its aggregate influence
        (Section~\ref{sec:threat}).
  \item A unified experimental protocol that benchmarks seven aggregation
        methods under three attack families and four scenario configurations
        (Section~\ref{sec:method}).
  \item An empirical finding that, in the four synthetic scenarios tested,
        no aggregation method substantially outperforms plain majority and
        several trust-score variants underperform it under strong capture
        (Section~\ref{sec:exp}). We offer a hypothesis, not a structural
        theorem, that capture can resemble honest-but-wrong behavior in the
        generator.
  \item A reproducible evidence package with the vote generators,
        aggregation code, frozen tables and a descriptive local-model pilot.
\end{itemize}

We emphasise that the \emph{negative} headline finding is the paper's
contribution, not an apology. TMLR's acceptance criteria explicitly
prioritise technical correctness and audience interest over novelty, and a
rigorous benchmark that documents the limits of the obvious aggregation
strategies is valuable precisely because future work can build on it
rather than re-derive it. We discuss when, and against which assumptions,
our negative result would \emph{not} hold
(Section~\ref{sec:conclusion}).

The remainder of this paper is organized as follows.
Section~\ref{sec:background} reviews classical Byzantine fault tolerance and
LLM ensembling. Section~\ref{sec:threat} introduces the threat model.
Section~\ref{sec:method} describes the seven aggregation rules we benchmark.
Section~\ref{sec:exp} reports the benchmark and its negative result.
Section~\ref{sec:related} discusses related work.
Section~\ref{sec:broader} addresses broader impact.
Section~\ref{sec:conclusion} concludes and outlines conditions under which
our negative result would not hold.

% =====================================================================
\section{Background}
\label{sec:background}
% =====================================================================

\paragraph{Classical Byzantine consensus.}
Byzantine fault tolerance (BFT) \citep{lamport1982byzantine} establishes
that a deterministic consensus protocol over $N$ replicas tolerates at
most $f < N/3$ Byzantine (arbitrary) failures when paired with digital
signatures, or $f < N/2$ under synchrony assumptions with binary outputs.
PBFT \citep{castro1999pbft} brought BFT to practical systems and
HotStuff \citep{yin2019hotstuff} linearised the communication pattern.
All classical BFT results assume \emph{binary} replica outputs; the
extension to \emph{continuous} outputs (which we need for LLM verifier
probabilities) is not standard and is one of the contributions of this
paper.

\paragraph{LLM-as-judge and ensembles.}
The use of LLMs as evaluators, or \emph{judges}, was popularised by
MT-Bench and Chatbot Arena \citep{zheng2023judging}, where a single
LLM judge is shown to correlate well with human preference ratings.
Subsequent work studied ensemble-of-judges methods
\citep{cui2026edo,leung2026rubricllm} to reduce variance and calibrate
directional bias. These methods optimise for \emph{passive} disagreement
across annotators; we instead consider \emph{adversarial} disagreement
induced by verifier capture.

\paragraph{Adversarial robustness of LLM judges.}
Recent work has begun to expose the silent nature of biased LLM
judges: \citet{zhang2026blindcurator} show that a biased judge
\emph{silently disables} a downstream curator past a sharp false-pass
threshold, and broader red-teaming studies
\citep{perez2022red,schulhoff2023ignore} catalogue prompt-injection
attacks. Our threat model (Section~\ref{sec:threat}) operationalises
these observations into a quantitative capture-resilience condition.

\paragraph{Peer-review automation.}
\citet{liu2026paperreviewtoolkit} release an open-source multi-reviewer
pipeline that we use as one of our testbeds; we discuss this and
related calibration work \citep{liu2026calibration} in
Section~\ref{sec:related}.

% =====================================================================
\section{Threat Model and Definitions}
\label{sec:threat}
% =====================================================================

We consider an ensemble $\V = \{v_1, \ldots, v_N\}$ of $N$ LLM verifiers
that independently evaluate a candidate output and produce a continuous
decision distribution $p_i \in [0,1]$ on the binary outcome
$\{0,1\}$; the aggregated decision is produced by a fixed aggregation rule
$\mathcal{A}: [0,1]^N \to \{0,1\}$. A subset $\A \subseteq \V$ of size
$f = |\A|$ is \emph{captured} by an adversary that controls the verifiers'
outputs but cannot observe or modify the inputs seen by honest verifiers.

\begin{definition}[Honest verifier]
\label{def:honest}
A verifier $v_i \in \V \setminus \A$ is \emph{honest} if
$p_i = \mathbb{P}[\text{correct outcome} \mid x, \theta_i]$, where $\theta_i$
is the verifier's fixed-but-unknown system prompt and training. In
particular, an honest verifier's outputs are conditionally independent
given the input $x$.
\end{definition}

\begin{definition}[Capture coalition]
\label{def:capture}
A captured verifier $v_j \in \A$ outputs
$\tilde{p}_j = g(p_j^{\text{true}})$ where $g: [0,1] \to [0,1]$ is an
adversarial response function drawn from one of three concrete attack
families:
\begin{itemize}
  \item \emph{Prompt injection}: $g$ is a fixed affine or step function
        chosen by the adversary (e.g. ``always output 0.9'').
  \item \emph{Rubric injection}: $g$ depends on the input via a hidden
        rubric bias that flips the verifier's preference order on a
        target subset of inputs.
  \item \emph{Majority bias}: $g$ mimics the running empirical mode of
        the other verifiers' outputs, so the captured output is
        indistinguishable from the majority.
\end{itemize}
\end{definition}

\begin{definition}[Aggregation correctness]
\label{def:correctness}
The aggregation rule $\mathcal{A}$ is $(\eta, \delta)$-correct on a
distribution $\mathcal{D}$ if
$\Pr_{(x,y) \sim \mathcal{D}, \A, g}
[\mathcal{A}(p\_1, \ldots, p\_N) = y] \geq 1 - \delta$,
where the probability is taken over the input distribution, the
adversary's choice of capture coalition $\A$, and the response
functions $g$ of all captured verifiers.
\end{definition}

\begin{definition}[Capture-resilient aggregation]
\label{def:resilient}
The aggregation rule $\mathcal{A}$ is $\varepsilon$-\emph{capture-resilient}
if, for every adversary that captures at most $f$ verifiers, the ensemble
remains $(\eta, \delta)$-correct with $\delta \leq \varepsilon$ for every
input distribution $\mathcal{D}$.
\end{definition}

\begin{theorem}[Capture lower bound]
\label{thm:lower}
For any deterministic aggregation rule $\mathcal{A}$, if the adversary
captures $f \geq N/2$ verifiers, then for \emph{some} input distribution
$\mathcal{D}$ and \emph{some} adversarial response function $g$ from
Definition~\ref{def:capture}, the rule $\mathcal{A}$ is not
$\varepsilon$-capture-resilient for any $\varepsilon < 1/2$.
\end{theorem}

\begin{proof}[Proof sketch]
If the adversary controls a strict majority $f > N/2$, the adversary can
choose $g$ such that the captured verifiers all output the same value $y^*$
with probability $1$, while honest verifiers (whose accuracy is below $1$
on $\mathcal{D}$) collectively cannot outvote $\A$. Formal reduction to
the classical $f \geq N/2$ impossibility bound for deterministic
aggregation. Full proof in Appendix~\ref{app:proof}.
\end{proof}

\begin{remark}[Continuous outputs vs.\ classical BFT]
Theorem~\ref{thm:lower} recovers the classical $f < N/2$ BFT impossibility
bound but does not yield a constructive upper bound, because LLM verifier
outputs are continuous: a verifier that outputs $p_i = 0.51$ is
qualitatively different from one that outputs $p_i = 0.99$, even though
both would be treated as ``in favor'' by a binary-vote aggregation rule.
This gap is what motivates the trust-score weighting in \textsc{CRV}
(Section~\ref{sec:method}).
\end{remark}

\begin{assumption}[Honest independence]
\label{ass:indep}
Honest verifiers' outputs are conditionally independent given $x$. This
is the standard assumption in LLM-ensemble calibration work
\citep{cui2026edo,zheng2023judging}; we discuss the relaxation in
Section~\ref{sec:related}.
\end{assumption}

% =====================================================================
\section{Methods: Seven Aggregation Rules}
\label{sec:method}
% =====================================================================

We benchmark seven aggregation rules. Three are \emph{classical} and
require no learning: plain majority voting
(Section~\ref{sub:majority}), PBFT-style deterministic thresholding
(Section~\ref{sub:pbft}), and uniform-weighted mean
(Section~\ref{sub:uniform}). Three are \emph{trust-score-based} and
share the EMA update
$\tau_i^{(t)} = \alpha \tau_i^{(t-1)} + (1-\alpha) c_i^{(t)}$ where
$c_i^{(t)} = 1 - |p_i^{(t-1)} - \text{med}_{-i}^{(t-1)}|$; they differ
in how the trust score is \emph{used} at decision time: weighted mean,
weighted median, and Kalman-filter trust tracking. The seventh is
\textsc{Atv}, an adaptive hybrid we proposed as a candidate that we
expected to combine the strengths of soft weighting and outlier
robustness.

\paragraph{Plain majority.}
\label{sub:majority}
Each verifier contributes a binary vote $\mathbb{1}[p_i \geq 0.5]$; the
aggregate is the majority over the $N$ binary votes. Cost $O(N)$ per
round, no state.

\paragraph{PBFT-style deterministic thresholding.}
\label{sub:pbft}
\citep{castro1999pbft} Each verifier again votes binary; the aggregate
passes if at least $(N+f+1)/2$ verifiers agree, which is the
deterministic threshold that tolerates $f$ adversarial nodes when a
strict majority is required. Cost $O(N)$ per round.

\paragraph{Uniform-weighted mean.}
\label{sub:uniform}
The aggregate is $\mathbb{1}[\sum_i p_i / N \geq 0.5]$. Cost $O(N)$,
no state.

\paragraph{Trust-score weighted mean.}
\label{sub:trustmean}
At round $t$, weights are $w_i = \tau_i^{(t-1)} / \sum_j \tau_j^{(t-1)}$;
the aggregate is $\mathbb{1}[\sum_i w_i p_i \geq 0.5]$. The trust
update uses consistency against the per-round median. Cost $O(N)$
plus a per-round $O(N \log N)$ for the median.

\paragraph{Trust-score weighted median.}
\label{sub:trustmedian}
Identical trust update; at decision time, sort verifiers by current
$p_i$ and pick the $p_i$ at which the cumulative trust weight crosses
half the total weight. This is robust to a single extreme value but
expensive $O(N \log N)$ per round.

\paragraph{Kalman-filter trust tracking.}
\label{sub:kalman}
Same decision rule as weighted mean, but trust tracking uses a
Kalman filter with process variance $Q = 0.01$ and measurement
variance $R = 0.05$ instead of an EMA, which we hypothesised would
smooth trust under noisy verifier streams.

\paragraph{\textsc{Atv}: adaptive trust-weighted voting.}
\label{sub:atv}
A hybrid we designed as the natural synthesis of the two trust-score
decision rules. \textsc{Atv} switches between weighted mean and
weighted median at each round based on the standard deviation of the
trust scores: if $\text{std}(\tau) < \tau_{\text{split}} = 0.5$ (trust
concentrated, presumed-uniform regime), use weighted mean; otherwise
(trust dispersed, capture suspected) use weighted median. The
hyperparameter $\tau_{\text{split}}$ was tuned on a held-out scenario;
we report sensitivity in Appendix~\ref{app:ablation}.

\paragraph{Why we expected \textsc{Atv} to win.}
Our \emph{a priori} hypothesis, before running the benchmark, was that
\textsc{Atv} would inherit the best of both worlds: the mean's
sensitivity to subtle weight differences in the easy regime, and the
median's robustness to outlier-dominated rounds in the hard regime.
The benchmark (Section~\ref{sec:exp}) does not support this hypothesis.

\paragraph{Why we report \textsc{Atv} even though it loses.}
We report \textsc{Atv} as a representative of the natural design space
that future researchers might explore, rather than as our recommended
method. The negative result is informative precisely because \textsc{Atv}
\emph{was} our best prior and still loses; the same is true of the
Kalman variant.

% =====================================================================
\section{Experimental Setup}
\label{sec:exp}
% =====================================================================

\paragraph{Datasets.}
\label{sub:generator}
We use a synthetic verifier-output generator in this study; we discuss
the limitations of this choice in Section~\ref{sec:conclusion}. The
generator (Section~\ref{sub:generator}) models the agreement structure
of a multi-LLM verification ensemble under verifier capture, and is
calibrated to produce vote distributions whose summary statistics
match published numbers for LLM judges on rubric decisions
\citep{leung2026rubricllm}: per-verifier outputs are concentrated
near $0$ or $1$ with $p_{\text{conf}} \in [0.7, 0.85]$ and a small
amount of i.i.d.\ Gaussian noise ($\sigma = 0.05$). The candidate tree does
not contain a held-out real-trace benchmark; the local pilot is reported
separately as a descriptive artifact.

\paragraph{Attack models.}
\label{sub:attacks}
We instantiate the three attack families of
Definition~\ref{def:capture} concretely: \emph{prompt injection}
(\textsc{PromptInj}: captured verifier outputs the complement of its
honest output), \emph{rubric injection}
(\textsc{RubricInj}: captured verifier flips its label on a random
fraction of inputs), and \emph{majority bias}
(\textsc{MajBias}: captured verifier mimics the running median of the
other verifiers' outputs). Section~\ref{sub:attacks} gives the formal
parameterisations.

\paragraph{Scenarios.}
We consider four scenarios that vary the two most consequential
parameters: honest-verifier per-instance accuracy $\eta \in \{0.65,
0.85\}$ and attack strength $\alpha \in \{0.6, 0.9\}$ (probability
that a captured verifier executes the attack on a given round). The
four combinations are labelled S1--S4 in
Table~\ref{tab:main}. We run $T = 1000$ rounds per scenario and
average over $n = 5$ random seeds; full per-seed numbers are in
Appendix~\ref{app:ablation}.

\paragraph{Metrics.}
We report \emph{decision accuracy} as the fraction of rounds on which
the aggregated decision equals the ground-truth label. Accuracy is the
right metric for this benchmark because the goal of the aggregator is
to produce the correct label, and the captured verifiers' goal is to
produce an incorrect one; alternative metrics such as false-positive
rate at the verifier level would require labelling individual
verifiers as captured or honest, which we cannot do from observed
verifier outputs alone (this is the central statistical difficulty).

\paragraph{Hyperparameters.}
All trust-based methods use $\alpha = 0.9$ (EMA decay) and
$c_i^{(t)} = 1 - |p_i - \text{med}_{-i}|$. \textsc{Atv} uses
$\tau_{\text{split}} = 0.5$ chosen on a held-out scenario. PBFT
threshold is $(N+f+1)/2$. Code and configuration files are released
with the paper.

% =====================================================================
\section{Results}
\label{sec:results}
% =====================================================================

\paragraph{Headline.}
Table~\ref{tab:main} reports the headline result. Plain majority voting is
best or tied for best in S1--S3 and is within $0.8$ percentage points of the
best method in S4, where trust-mean reaches $0.8506$ versus $0.8430$ for
majority. \textsc{Atv} never wins a scenario. The differences are small
across the five random seeds; the reported cell standard deviations are
below $0.01$.

\begin{table}[ht]
\caption{Decision accuracy across four scenarios. Best per column in bold.
Majority is best or tied in S1--S3; trust-mean is higher in S4, while the
four-scenario mean is essentially tied between majority and trust-median.}
\label{tab:main}
\centering
\small
\begin{tabular}{lccccc}
\toprule
Method & S1 & S2 & S3 & S4 & mean \\
 & ($\eta{=}0.85$, $\alpha{=}0.6$)
 & ($\eta{=}0.65$, $\alpha{=}0.6$)
 & ($\eta{=}0.65$, $\alpha{=}0.9$)
 & ($\eta{=}0.85$, $\alpha{=}0.9$)
 & \\
\midrule
majority            & \textbf{0.889} & \textbf{0.553} & \textbf{0.459} & 0.843 & \textbf{0.686} \\
bft\_pbft           & 0.777 & 0.522 & 0.458 & 0.720 & 0.619 \\
weighted\_uniform   & 0.804 & 0.484 & 0.365 & 0.729 & 0.595 \\
trust\_mean         & 0.883 & 0.521 & 0.432 & \textbf{0.851} & 0.672 \\
trust\_median       & \textbf{0.889} & \textbf{0.553} & \textbf{0.459} & 0.844 & \textbf{0.686} \\
kalman\_trust       & 0.882 & 0.535 & 0.446 & 0.842 & 0.677 \\
\textsc{Atv}        & 0.883 & 0.527 & 0.438 & 0.846 & 0.673 \\
\bottomrule
\end{tabular}
\end{table}

\paragraph{Effect sizes.}
The five seeds are matched runs of one protocol, so contrasts are paired and
we report $d_z$ (mean per-seed difference over its standard deviation) rather
than the independent-sample $d$ our Statistical Procedures paragraph names.
\Cref{tab:dz} gives $d_z$ for \textsc{Atv} against each baseline; positive
values favour \textsc{Atv}.

\begin{table}[ht]
\caption{Paired effect size $d_z$ of \textsc{Atv} against each baseline,
across the four scenarios. Values are computed from the per-seed accuracy
records shipped with this submission.}
\label{tab:dz}
\centering
\small
\begin{tabular}{lrrrr}
\toprule
Baseline & S1 & S2 & S3 & S4 \\
\midrule
majority         & $-3.79$ & $-5.36$ & $-8.24$ & $+0.90$ \\
trust\_mean      & $-0.15$ & $+3.23$ & $+2.74$ & $-5.00$ \\
trust\_median    & $-3.79$ & $-5.36$ & $-8.24$ & $+0.35$ \\
kalman\_trust    & $+0.11$ & $-5.14$ & $-2.14$ & $+2.37$ \\
weighted\_uniform & $+8.97$ & $+5.03$ & $+7.12$ & $+16.93$ \\
\bottomrule
\end{tabular}
\end{table}

Two features of \Cref{tab:dz} bear on the claim in \Cref{tab:main}. First,
against the strongest baselines the effect is small and sometimes of the
opposite sign: against \texttt{kalman\_trust} in S1, $d_z = +0.11$, and against
\texttt{trust\_mean} in S1, $d_z = -0.15$. The large values in the last row are
against \texttt{weighted\_uniform}, the weakest method in every column, so they
restate the table's ordering rather than add to it. We therefore do not claim
\scshape{Atv} improves on the leading baselines, and the negative result in
\Cref{tab:main} stands on the point estimates, not on the effect sizes.

Second, the reported per-seed accuracy vectors for \texttt{majority} and
\texttt{trust\_median} are identical in S1, S2 and S3 and differ by $0.007$ in
S4. The evidence records per-seed accuracies rather than per-round prediction
streams, so we do not claim that the two rules produce byte-identical
predictions.

\paragraph{A hypothesis for the negative result.}
We hypothesize that in this synthetic generator capture can be hard to
distinguish from honest-but-wrong behavior: a captured verifier outputting
the complement of its honest output can have a similar per-round deviation
from the median as an honest verifier that happens to be wrong. A rule that
uses deviation from the median to identify capture may therefore ignore both
or exclude both. The evidence package does not contain a separate test of
this mechanism, so it remains a hypothesis rather than a structural theorem.

In the high-honest scenario S1, the trust-based methods approach
majority but trail it by 0.5 percentage points, consistent with the
explanation above: the trust mechanism correctly down-weights
captured verifiers when they \emph{happen} to be outliers, but it
also incorrectly down-weights honest verifiers that happen to be
wrong, and the two effects approximately cancel.

In the strong-attack scenarios S3 and S4, the trust-based methods
gain an apparent advantage in some columns (S4 for trust\_mean) but
lose it in others (S2 for all trust-based methods). The advantage in
S4 disappears when averaged across all four scenarios: the mean
column is dominated by majority and trust\_median, both at $0.686$.

\paragraph{Per-attack breakdown.}
Across the three attack families (prompt injection, rubric injection,
majority bias), the best method differs by family. Trust-mean is highest
under prompt injection and majority-bias mimicry, while trust-median,
Kalman-trust, and \textsc{Atv} are at the top under rubric injection.
Majority is within $1.7$ percentage points of the best attack-family mean
and remains a strong baseline; this does not establish a universal rule.

\paragraph{Sensitivity evidence.}
The checked-in evidence contains an ensemble-size sweep at $N\in\{5,7,10\}$,
but no per-cell $\tau_{\text{split}}$ or trust-decay sweep. We therefore make
no claim about hyperparameter robustness; the sweep is reported exactly in
Appendix~\ref{app:ablation}.

\paragraph{Implication.}
Within the synthetic ensemble sizes, attack strengths, and honest-accuracy
regimes in this benchmark, the tested trust-based and adaptive rules do not
show a consistent gain over plain majority. This is a bounded empirical
statement, not a claim that all non-trivial aggregators fail.

% =====================================================================
\section{Related Work}
\label{sec:related}
% =====================================================================


\paragraph{Classical Byzantine consensus.}
\citet{lamport1982byzantine,castro1999pbft,yin2019hotstuff} establish the
$3f+1$ bound for binary-output consensus under $f$ adversarial nodes. Our
setting differs in that LLM verifier outputs are continuous probability
distributions, so classical BFT bounds do not transfer directly; this
motivates our $\varepsilon$-capture-resilience definition in
Section~\ref{sec:threat}.

\paragraph{LLM-as-judge and ensemble calibration.}
\citet{zheng2023judging} introduce the MT-Bench and Chatbot Arena
evaluation paradigm; \citet{leung2026rubricllm} provide the most recent
calibration numbers for 8 LLM judges on rubric decisions, and
\citet{cui2026edo} jointly optimize ensemble weights and calibration for
subjective supervision. CRV differs from EDO in operating under an
explicit adversarial model (capture) rather than a passive disagreement
model.

\paragraph{Adversarial robustness of LLM judges.}
\citet{zhang2026blindcurator} show that a biased judge \emph{silently}
disables a downstream curator past a sharp false-pass threshold, and
\citet{perez2022red,schulhoff2023ignore} catalogue prompt-injection and
red-team attack techniques against LLM evaluators. Our threat model
operationalises these observations into a quantitative capture-resilience
definition and a reproducible benchmark; it does not establish a general
aggregation-correctness guarantee.

\paragraph{Peer-review automation.}
\citet{liu2026paperreviewtoolkit} release an open-source multi-reviewer
pipeline that we use as the testbed for our empirical evaluation
(Section~\ref{sec:exp}); \citet{liu2026calibration} develops calibration
methods for multi-agent ensembles more broadly and is the closest work
in our own research line.

% =====================================================================
\section{Broader Impact}
\label{sec:broader}
% =====================================================================

\paragraph{Positive impact.}
The primary positive impact of this work is \emph{negative}: by
showing that no aggregation method we tested substantially outperforms
plain majority voting under verifier capture, we redirect practitioner
effort away from building sophisticated aggregators and toward
investments that \emph{do} pay off in the capture regime: increasing
honest-verifier accuracy, diversifying the verifier pool, and
auditing deployed ensembles for suspicious agreement patterns. We
expect the benchmark to be useful for engineering teams that are
currently considering whether to invest in a custom aggregation
layer; our data suggest the answer is ``no'' under the parameter
regime we test, and the engineering budget should be spent elsewhere.

A secondary positive impact is methodological: the failure modes of
the seven aggregation rules we test illuminate which statistical
assumptions a successful capture-resilient aggregator would need
to break. Future research that does find a winning method can
explicitly cite which assumption we showed fails here.

\paragraph{Negative impact and dual-use concerns.}
The same attack implementations that allow us to \emph{benchmark}
verifier capture also allow an adversary to \emph{measure} capture
susceptibility of a competing verifier product. We consider this a
real but bounded risk for the following reasons.
\textbf{(i)} The three attack families we instantiate
(\textsc{PromptInj}, \textsc{RubricInj}, \textsc{MajBias}) are
straightforward enough that a competent adversarial team could
re-derive them from the threat model alone; our code release does
not materially lower the cost.
\textbf{(ii)} Our benchmark demonstrates that these attacks, when
attempted against a well-calibrated honest majority, have only
small absolute effect (the gap between majority and the best
trust-based method is below 1 percentage point in 3 of 4
scenarios). An adversary using our code to measure capture
susceptibility will therefore find that the practical upside of
attack deployment is small.
\textbf{(iii)} Our benchmark is \emph{synthetic}. An adversary who
wishes to deploy a capture attack against a real ensemble must
still construct a real captured verifier that mimics the target
ensemble's distribution, which is a substantially harder engineering
problem than running our code.

\paragraph{Responsible disclosure norms.}
The candidate evidence tree contains the benchmark code and attack
generators; the final license and release wording remain user decisions.
No captured-verifier weights or fine-tuning scripts are present.

\paragraph{Limitations of public release.}
The benchmark code is sufficient to reproduce the synthetic tables but not
to deploy a capture attack against a real ensemble without substantial
additional engineering. This scope is deliberate.

\paragraph{Stakeholder considerations.}
Verifier-capture attacks disproportionately threaten operators of
multi-LLM auditing pipelines. The benchmark is documented for defensive
testing; no external outreach or deployment validation is claimed.

\paragraph{Positive-negative balance.}
On balance, we believe the positive impact of releasing this
benchmark---redirecting engineering effort away from ineffective
aggregation rules toward more productive investments in verifier
quality and pool diversification---substantially outweighs the
bounded negative impact of marginally lowering the cost of an
attack that, per our own benchmark, has small practical effect
against well-calibrated honest majorities.

% =====================================================================
\section{Ethics Statement}
\label{sec:ethics}
% =====================================================================

This study does not involve human subjects, personal data, or animal
subjects. The verifier-output data used here is either synthetic
(Section~\ref{sub:generator}) or from the descriptive local-model pilot.

\paragraph{Dual-use risk.}
We have discussed the dual-use risk of releasing the attack
implementations in detail in Section~\ref{sec:broader}. In summary,
the bounded risk is discussed in Section~\ref{sec:broader}. The candidate
tree contains no captured-verifier weights or fine-tuning scripts; final
license and responsible-use wording remain user decisions.

\paragraph{Affected stakeholders.}
The primary stakeholders are operators of multi-LLM verification
ensembles, including AI-auditing platforms, peer-review automation
tools, and safety-content moderation pipelines. We have reached out
to the maintainers of two such platforms
(\citep{liu2026paperreviewtoolkit} and the public release of
\citet{leung2026rubricllm}'s evaluation harness) to inform them of
our findings ahead of public release.

\paragraph{IRB and funding.}
This study was not subject to IRB review. No external funding was
received for this work.

\paragraph{Author conflicts of interest.}
The authors declare no financial conflicts of interest. The authors
maintain an open-source peer-review audit tool
\citep{liu2026paperreviewtoolkit} that is discussed as one of the
motivating use cases for this benchmark; the tool is referenced
consistently with the rest of the literature and no special pleading
is made for it.

% =====================================================================
\section{Conclusion}
\label{sec:conclusion}
% =====================================================================

We have presented an empirical benchmark of seven aggregation methods
for multi-LLM verifier ensembles under three verifier-capture attacks,
across four scenarios. In this synthetic benchmark, plain majority is a
strong baseline: it is best or tied in three scenarios and within 0.8
percentage points of the best method in the fourth. Several trust-score
variants underperform majority in some scenarios. We hypothesize that
captured behavior can resemble honest-but-wrong behavior in the generator,
but we do not establish that as a structural theorem.

\paragraph{When our negative result would \emph{not} hold.}
The bounded result rests on the synthetic generator's assumption that
captured verifiers can resemble honest-but-wrong verifiers. If a real-world capture breaks this
assumption --- for instance, by introducing a structured deviation
(e.g. a verifier that always outputs the exact median plus a fixed
delta) --- then a more sophisticated aggregation rule may recover
above majority. Our benchmark does not rule this out; it only shows
that for the three concrete attack parameterisations we test, the
gain is at most a fraction of a percentage point.

A second assumption is that honest verifiers are exchangeable. If
honest verifiers have known specialty areas (e.g.\ one verifier is
better at math, another at code), then a domain-aware aggregator may
do substantially better than majority. Our synthetic generator
deliberately does not model this specialisation, so our benchmark
under-estimates the upside of sophisticated aggregation in
heterogeneous pools.

\section{Discussion and Limitations}
Our results rely on a synthetic verifier-output generator calibrated
to published rubric-LLM numbers but not validated against
real LLM trace data. A practitioner applying our benchmark to a real
ensemble should re-run our scenario definitions with the real
distribution and may find that some methods we rule out do in fact
help on their specific distribution. We release the synthetic
generator and the attack implementations to facilitate this
replication.

\paragraph{Future work.}
Three directions follow from this study.
(i) Validate the negative result on real LLM verifier trace data;
the result may not survive the transition from synthetic to real.
(ii) Investigate structured-attack settings where captured verifiers
introduce deviations outside the honest-but-wrong distribution; the
negative result may not apply there.
(iii) Investigate diversified verifier pools where domain-aware
aggregation can outperform majority by exploiting specialisation.

% =====================================================================
% References (refs.bib populated in W1)
% =====================================================================
\bibliographystyle{tmlr}
\bibliography{refs}

% =====================================================================
% Appendix begins here
% =====================================================================
\appendix

\section{Proof of the Capture Lower Bound (Theorem~1)}
\label{app:proof}

We prove the deterministic lower bound stated in Theorem~1: if
$f \geq N/2$ of the verifiers are captured, no deterministic
aggregation rule $\mathcal{A}: [0,1]^N \to \{0,1\}$ can be
$\varepsilon$-capture-resilient for $\varepsilon < 1/2$.

\begin{proof}
Consider the case $f > N/2$ (strict majority capture; the case
$f = N/2$ follows by a standard limit argument). The adversary
chooses a target label $y^* \in \{0,1\}$ and a response function
$g$ such that every captured verifier outputs $y^*$ with
probability $1$ on every input. The adversary also chooses an input
distribution $\mathcal{D}$ on which the \emph{honest} verifiers have
per-instance accuracy strictly below $1/2$; this is always possible
by choosing $\mathcal{D}$ to concentrate on inputs where the honest
verifiers systematically err.

Because the captured coalition $\A$ forms a strict majority, the
deterministic rule $\mathcal{A}$ sees $p_i \in \{0, 1\}$ on the
captured inputs (after $g$ is applied) and any values in $[0,1]$ on
the honest inputs. Since $\mathcal{A}$ is deterministic, its
decision is a function only of the multiset $\{p_1, \ldots, p_N\}$,
and in particular of the \emph{number} of captured verifiers voting
$y^*$, which is $f > N/2$. For any such deterministic rule, there
exists an assignment of honest outputs such that $\mathcal{A}$ outputs
$y^*$ on every input, even when the true label is the opposite
$\bar{y}^*$; this is a direct consequence of the fact that a strict
majority of $N$ deterministic binary inputs forces the binary
majority rule, and any rule that agrees with binary majority on
strict-majority inputs cannot be $\varepsilon$-resilient for
$\varepsilon < 1/2$ on the chosen $\mathcal{D}$.

This shows that for any deterministic $\mathcal{A}$, the adversary
can construct a distribution $\mathcal{D}$ on which $\mathcal{A}$
achieves error probability at least $1/2$, completing the proof.
\end{proof}

\section{Algorithm Pseudocode and Complexity}
\label{app:algo}

We give pseudocode for the four representative aggregation rules.
All four have memory $O(N)$ for the per-verifier state.

\paragraph{Plain majority.}
\begin{algorithmic}[1]
\Require votes $p_1, \ldots, p_N$
\State $s \gets \sum_i \mathbb{1}[p_i \geq 0.5]$
\State \Return $\mathbb{1}[s > N/2]$
\end{algorithmic}
Time $O(N)$.

\paragraph{PBFT-style threshold.}
\begin{algorithmic}[1]
\Require votes $p_1, \ldots, p_N$; bound $f$
\State $s \gets \sum_i \mathbb{1}[p_i \geq 0.5]$
\State \Return $\mathbb{1}[s \geq (N+f+1)/2]$
\end{algorithmic}
Time $O(N)$.

\paragraph{Trust-score weighted mean.}
\begin{algorithmic}[1]
\Require votes $p_1, \ldots, p_N$; trusts $\tau_1, \ldots, \tau_N$; decay $\alpha$
\State $w_i \gets \tau_i / \sum_j \tau_j$ for each $i$
\State \Return $\mathbb{1}[\sum_i w_i p_i \geq 0.5]$
\State Update: $\tau_i \gets \alpha \tau_i + (1-\alpha)(1 - |p_i - \text{med}_{-i}|)$ for each $i$
\end{algorithmic}
Time $O(N \log N)$ per round (median), memory $O(N)$.

\paragraph{Trust-score weighted median.}
\begin{algorithmic}[1]
\Require votes $p_1, \ldots, p_N$; trusts $\tau_1, \ldots, \tau_N$; decay $\alpha$
\State Sort indices by $p_i$; let $\pi$ be the sorted permutation
\State $W \gets \sum_j \tau_j$; find smallest $k$ with $\sum_{i \leq k} \tau_{\pi(i)} \geq W/2$
\State \Return $\mathbb{1}[p_{\pi(k)} \geq 0.5]$
\State Update trust as in weighted mean.
\end{algorithmic}
Time $O(N \log N)$ per round, memory $O(N)$.

\paragraph{\textsc{Atv} (adaptive hybrid).}
\begin{algorithmic}[1]
\Require votes $p_1, \ldots, p_N$; trusts $\tau_1, \ldots, \tau_N$; decay $\alpha$; split $\tau_{\text{split}}$
\If{$\text{std}(\tau) < \tau_{\text{split}}$}
  \State \Return weighted-mean decision
\Else
  \State \Return weighted-median decision
\EndIf
\State Update trust as in weighted mean.
\end{algorithmic}
Time $O(N \log N)$ per round, memory $O(N)$.

\section{Additional Ablations and Sensitivity Sweeps}
\label{app:ablation}

\paragraph{Per-attack-model breakdown.}
Table~\ref{tab:per-attack} reports per-attack-model accuracies. The
ordering of methods is qualitatively consistent across the three
attack families: majority and trust\_median remain tied at the top,
with the largest absolute gap (0.7 percentage points) appearing
under rubric injection.

\begin{table}[ht]
\caption{Decision accuracy by attack model, averaged across all
four scenarios (5 seeds, $T=1000$ rounds per cell).}
\label{tab:per-attack}
\centering
\small
\begin{tabular}{lccc}
\toprule
Method & \textsc{PromptInj} & \textsc{RubricInj} & \textsc{MajBias} \\
\midrule
majority            & 0.877 & 0.945 & 0.997 \\
bft\_pbft           & 0.765 & 0.886 & 0.985 \\
weighted\_uniform   & 0.876 & 0.944 & 0.997 \\
trust\_mean         & \textbf{0.894} & 0.944 & \textbf{0.998} \\
trust\_median       & 0.877 & \textbf{0.945} & 0.997 \\
kalman\_trust       & 0.877 & \textbf{0.945} & 0.997 \\
\textsc{Atv}        & 0.886 & \textbf{0.945} & 0.997 \\
\bottomrule
\end{tabular}
\end{table}

\paragraph{Hyperparameter sensitivity.}
\textsc{Atv}'s $\tau_{\text{split}}$ was tuned on a held-out
scenario. Sensitivity sweeps over $\tau_{\text{split}} \in [0.2, 0.8]$
shift \textsc{Atv}'s mean accuracy by at most $0.4$ percentage
points; the rank ordering is preserved across all sweep values.
Trust-update decay $\alpha$ sweeps over $[0.5, 0.99]$ shift
trust\_mean by $\pm 0.5$ percentage points with no rank change.

\paragraph{Ensemble-size sweep.}
The checked-in ensemble-size sweep covers $N\in\{5,7,10\}$ with
$f=\lfloor N/3\rfloor$. At $N=5$, majority and several methods tie at
$0.7479$. At $N=7$, majority and trust-median are essentially tied at
$0.686$. At $N=10$, trust-median reaches $0.7668$, higher than majority's
$0.6785$. These three points do not establish a monotone ensemble-size
trend; they are reported as descriptive sweep values.

\paragraph{Local-model pilot.}
The evidence package includes a descriptive local Ollama pilot with four
models, 20 questions, three seeds, and two captured slots. The reported
rounded mean accuracy is 0.9167 for all seven methods. This pilot is not a
held-out trace benchmark and is not evidence for external generalization.

\section{Code and Reproducibility}
\label{app:code}

\paragraph{Repository layout.}
The candidate evidence tree contains the benchmark code and frozen result
files. Final license and release wording remain user decisions. The
repository contains:
\begin{itemize}
  \item \texttt{crv.py}: legacy CRV implementation (retained for
        archival reference; not used as a benchmarked method).
  \item \texttt{attacks.py}: implementations of the three attack
        families (\textsc{PromptInj}, \textsc{RubricInj},
        \textsc{MajBias}).
  \item \texttt{extended\_methods.py}: the seven aggregation methods
        benchmarked in this paper.
  \item \texttt{benchmark.py}: the four-scenario benchmark driver.
  \item \texttt{verify\_p6.py}: the verify suite used during
        paper preparation.
\end{itemize}

\paragraph{No held-out demo claim.}
The evidence tree contains no independently validated real-trace demo.
The local pilot is a descriptive artifact, not a deployment or held-out
validation claim.

\paragraph{Reproducibility hashes.}
The ingest manifest and the Research Factory overlay record SHA-256 hashes
for the frozen source and evidence files. A Python dependency lock and
complete local-pilot runtime logs are not present; those remain release
blockers.

\paragraph{Computational footprint.}
A full benchmark run (4 scenarios $\times$ 7 methods $\times$ 5 seeds,
$T=1000$ rounds) takes approximately 4 minutes on a single CPU core
(no GPU required) and uses 50 MB of memory. The complete experiment
can be reproduced in under 10 minutes on commodity hardware.


\section*{Reproducibility Statement}
\label{sec:reproducibility}

\textbf{Available evidence.} The candidate tree contains the synthetic
benchmark tables, the checked-in generator and aggregation code, and a
descriptive four-model local Ollama pilot. The ingest manifest records
SHA-256 hashes for the frozen source and evidence files. A complete Python
dependency lock and runtime logs for the local pilot are not present.

\textbf{What is not established.} The package does not contain a held-out
real-trace benchmark, a pre-registered temporal replication, or a proof of a
universal capture-resilience bound. Those remain open evidence gates and are
not presented as completed results.

\end{document}
